\documentclass[12pt,a4paper]{article}
\usepackage{mathwizards-ingles}
\title{%
  \vspace{-1cm}
  {\Huge\bfseries\textcolor{mes3}{Mes 3 · Semanas 9--10}}\\[0.3cm]
  {\Large Comprensión Auditiva, Debate y Reported Speech}\\[0.2cm]
  {\large\color{grissuave} Nivel A2 $\rightarrow$ B1 · Clubes de Escucha y Producción Oral}\\[0.5cm]
  \rule{\textwidth}{1.5pt}
}
\author{}
\date{}

\begin{document}
\maketitle
\thispagestyle{empty}

\begin{cajanivel}
\textbf{Objetivo:} El estudiante comienza a consumir \textbf{podcasts} y participa en \textbf{debates guiados}. Se introduce el \textbf{Past Continuous} para narrativas, el \textbf{reported speech} básico y los \textbf{discourse markers} necesarios para articular opiniones en los clubes de escucha.
\end{cajanivel}

% ═══════════════════════════════════════════════════════════════════════════════
\section{Gramática Emergente: Past Continuous}
% ═══════════════════════════════════════════════════════════════════════════════

\begin{cajagrammar}[Past Continuous: \textit{was/were} + verbo + \textbf{-ing}]
Se usa para describir acciones \textbf{en progreso en un momento del pasado}, frecuentemente interrumpidas por otra acción en Past Simple.

\medskip
\begin{tabularx}{\textwidth}{@{} l X @{}}
\toprule
\textbf{Estructura} & \textbf{Ejemplo} \\
\midrule
\textbf{Afirmativo} & \textit{I \textbf{was walking} home at 8 p.m.} \\
\textbf{Negativo} & \textit{She \textbf{wasn't listening} to the teacher.} \\
\textbf{Pregunta} & \textit{\textbf{Were} you \textbf{sleeping} when I called?} \\
\bottomrule
\end{tabularx}
\end{cajagrammar}

\begin{cajagrammar}[Combinación Narrativa: Past Continuous + Past Simple]
Este es el patrón estrella de la narración: una acción larga (continuous) interrumpida por una acción corta (simple).

\medskip
\textbf{Estructura:} \eng{While} / \eng{When} + Past Continuous, Past Simple.

\begin{cajaejemplo}[Ejemplos Narrativos]
\begin{itemize}[nosep, leftmargin=0.8cm]
  \item \textit{I \textbf{was walking} home \textbf{when} it \textbf{started} to rain.}
  \item \textit{\textbf{While} she \textbf{was cooking}, the phone \textbf{rang}.}
  \item \textit{They \textbf{were playing} football \textbf{when} the teacher \textbf{arrived}.}
  \item \textit{\textbf{While} I \textbf{was reading}, I \textbf{heard} a loud noise.}
\end{itemize}
\end{cajaejemplo}

\medskip
\begin{cajaejemplo}[Diferencia con Past Simple solo]
\begin{itemize}[nosep, leftmargin=0.8cm]
  \item Past Simple + Past Simple = dos acciones secuenciales:\\
    \textit{I \textbf{opened} the door and \textbf{saw} a dog.} (primero abrí, luego vi)
  \item Past Continuous + Past Simple = una acción interrumpe a otra:\\
    \textit{I \textbf{was opening} the door when I \textbf{saw} a dog.} (estaba abriendo y vi)
\end{itemize}
\end{cajaejemplo}
\end{cajagrammar}

% ═══════════════════════════════════════════════════════════════════════════════
\section{Gramática Emergente: Reported Speech Básico}
% ═══════════════════════════════════════════════════════════════════════════════

\begin{cajagrammar}[Reported Speech — Estilo Indirecto Básico]
Se usa para \textbf{contar lo que alguien dijo} sin citar textualmente. Los tiempos verbales ``retroceden'' un paso.

\medskip
\begin{tabularx}{\textwidth}{@{} l l @{}}
\toprule
\textbf{Discurso Directo} & \textbf{Discurso Indirecto (Reported)} \\
\midrule
\textit{``I \textbf{am} tired.''} & \textit{He \textbf{said} (that) he \textbf{was} tired.} \\
\textit{``I \textbf{like} pizza.''} & \textit{She \textbf{said} (that) she \textbf{liked} pizza.} \\
\textit{``I \textbf{will} call you.''} & \textit{He \textbf{said} (that) he \textbf{would} call me.} \\
\textit{``I \textbf{can} help.''} & \textit{She \textbf{said} (that) she \textbf{could} help.} \\
\textit{``I \textbf{have finished}.''} & \textit{He \textbf{said} (that) he \textbf{had finished}.} \\
\bottomrule
\end{tabularx}

\medskip
\textbf{Cambios de referencia:}
\begin{multicols}{2}
\begin{tabularx}{\linewidth}{@{} l l @{}}
I $\rightarrow$ & he / she \\
my $\rightarrow$ & his / her \\
we $\rightarrow$ & they \\
\end{tabularx}

\columnbreak

\begin{tabularx}{\linewidth}{@{} l l @{}}
today $\rightarrow$ & that day \\
tomorrow $\rightarrow$ & the next day \\
here $\rightarrow$ & there \\
\end{tabularx}
\end{multicols}

\medskip
\begin{cajaejemplo}[Reported Speech con preguntas]
\textbf{Preguntas sí/no:} se introduce con \textbf{if / whether}:\\
\textit{``Do you like coffee?''} $\rightarrow$ \textit{She \textbf{asked} me \textbf{if} I liked coffee.}

\medskip
\textbf{WH-questions:} se mantiene la palabra interrogativa pero con orden afirmativo:\\
\textit{``Where do you live?''} $\rightarrow$ \textit{He \textbf{asked} me \textbf{where} I lived.}
\end{cajaejemplo}
\end{cajagrammar}

% ═══════════════════════════════════════════════════════════════════════════════
\section{Chunks: Discourse Markers para Debate}
% ═══════════════════════════════════════════════════════════════════════════════

\begin{cajachunk}[Discourse Markers — Conectores para Articular Opiniones]
Estos marcadores son esenciales para participar en los \textbf{Clubes de Escucha y Debate}.

\medskip
\textbf{Para dar tu opinión:}
\begin{tabularx}{\textwidth}{@{} X l @{}}
\eng{I think / believe that...} & \esp{Creo que...} \\
\eng{In my opinion, ...} & \esp{En mi opinión, ...} \\
\eng{From my point of view, ...} & \esp{Desde mi punto de vista, ...} \\
\eng{As far as I'm concerned, ...} & \esp{En lo que a mí respecta, ...} \\
\end{tabularx}

\medskip
\textbf{Para estar de acuerdo o en desacuerdo:}
\begin{tabularx}{\textwidth}{@{} X l @{}}
\eng{I agree with you.} & \esp{Estoy de acuerdo contigo.} \\
\eng{I agree because...} & \esp{Estoy de acuerdo porque...} \\
\eng{I see your point, but...} & \esp{Entiendo tu punto, pero...} \\
\eng{I'm not sure about that.} & \esp{No estoy seguro de eso.} \\
\eng{I disagree because...} & \esp{No estoy de acuerdo porque...} \\
\end{tabularx}

\medskip
\textbf{Para estructurar el argumento:}
\begin{tabularx}{\textwidth}{@{} X l @{}}
\eng{First of all, ...} & \esp{En primer lugar, ...} \\
\eng{Moreover, ...} / \eng{Furthermore, ...} & \esp{Además, ...} \\
\eng{However, ...} & \esp{Sin embargo, ...} \\
\eng{Although...} & \esp{Aunque...} \\
\eng{On the other hand, ...} & \esp{Por otro lado, ...} \\
\eng{In conclusion, ...} & \esp{En conclusión, ...} \\
\eng{To sum up, ...} & \esp{En resumen, ...} \\
\end{tabularx}
\end{cajachunk}

% ═══════════════════════════════════════════════════════════════════════════════
\section{Chunks: Hablar de Podcasts y Medios}
% ═══════════════════════════════════════════════════════════════════════════════

\begin{cajachunk}[Discussing Media — Frases para los Clubes de Escucha]
\begin{tabularx}{\textwidth}{@{} X l @{}}
\eng{The podcast was about...} & \esp{El podcast fue sobre...} \\
\eng{The main topic was...} & \esp{El tema principal fue...} \\
\eng{The speaker said that...} & \esp{El hablante dijo que...} \\
\eng{I found it interesting because...} & \esp{Me pareció interesante porque...} \\
\eng{The part I liked the most was...} & \esp{La parte que más me gustó fue...} \\
\eng{I didn't understand the part about...} & \esp{No entendí la parte sobre...} \\
\eng{It reminded me of...} & \esp{Me recordó a...} \\
\eng{I learned that...} & \esp{Aprendí que...} \\
\eng{What do you think about...?} & \esp{¿Qué opinas sobre...?} \\
\eng{Do you agree with what the speaker said?} & \esp{¿Estás de acuerdo con lo que dijo?} \\
\end{tabularx}
\end{cajachunk}

% ═══════════════════════════════════════════════════════════════════════════════
\section{Chunks: Conectar Ideas con \textit{because, but, so}}
% ═══════════════════════════════════════════════════════════════════════════════

\begin{cajachunk}[Advanced Connectors — Para Argumentar en los Debates]
\begin{tabularx}{\textwidth}{@{} l l X @{}}
\toprule
\textbf{Conector} & \textbf{Función} & \textbf{Ejemplo} \\
\midrule
\eng{because} & causa & \textit{I liked it \textbf{because} it was very clear.} \\
\eng{so} & consecuencia & \textit{I didn't understand, \textbf{so} I listened again.} \\
\eng{but} & contraste & \textit{The topic was boring, \textbf{but} the speaker was good.} \\
\eng{although} & contraste fuerte & \textit{\textbf{Although} it was long, I enjoyed it.} \\
\eng{however} & contraste formal & \textit{The story was simple. \textbf{However}, the ending was surprising.} \\
\eng{moreover} & adición formal & \textit{It was informative. \textbf{Moreover}, it was entertaining.} \\
\eng{therefore} & conclusión & \textit{She practiced every day. \textbf{Therefore}, she improved quickly.} \\
\bottomrule
\end{tabularx}
\end{cajachunk}

% ═══════════════════════════════════════════════════════════════════════════════
\section{Oraciones Listas para Minar}
% ═══════════════════════════════════════════════════════════════════════════════

\begin{cajamineria}[Oraciones \textit{i+1} — Semanas 9--10]

\begin{enumerate}[leftmargin=1.2cm]
  \item \textit{I \underline{was listening} to a podcast when my phone rang.}\\
    {\small\textbf{Objetivo:} Past Continuous interrumpido por Past Simple.}

  \item \textit{While we \underline{were having} dinner, the lights went out.}\\
    {\small\textbf{Objetivo:} estructura ``While + Past Continuous, Past Simple''.}

  \item \textit{She said that she \underline{would} come to the meeting.}\\
    {\small\textbf{Objetivo:} reported speech con retroceso ``will $\rightarrow$ would''.}

  \item \textit{He asked me \underline{if} I liked the new restaurant.}\\
    {\small\textbf{Objetivo:} reported speech para preguntas sí/no con ``if''.}

  \item \textit{The podcast was interesting. \underline{However}, it was a bit long.}\\
    {\small\textbf{Objetivo:} el discourse marker ``However'' para contrastar.}

  \item \textit{I agree with you \underline{because} the speaker made very good points.}\\
    {\small\textbf{Objetivo:} la estructura ``I agree because...'' para debatir.}

  \item \textit{\underline{Although} I didn't understand everything, I got the main idea.}\\
    {\small\textbf{Objetivo:} el conector ``Although'' al inicio de oración.}

  \item \textit{The teacher asked us \underline{where} we had traveled.}\\
    {\small\textbf{Objetivo:} reported speech de WH-question con orden afirmativo.}

  \item \textit{\underline{On the other hand}, some people think it's too difficult.}\\
    {\small\textbf{Objetivo:} la expresión ``On the other hand'' para presentar el otro lado.}

  \item \textit{She practiced every day. \underline{Therefore}, her pronunciation improved a lot.}\\
    {\small\textbf{Objetivo:} el conector ``Therefore'' para conclusiones.}
\end{enumerate}
\end{cajamineria}

\vfill
\begin{center}
\rule{0.5\textwidth}{0.5pt}\\[0.3cm]
{\small\color{grissuave} Curso de Inmersión en Inglés · Mes 3 · Semanas 9--10 · Nivel A2 $\rightarrow$ B1}
\end{center}

\end{document}
